% !TEX program = lualatex
\documentclass[11pt,a4paper]{article}
\usepackage{luatexja}
\usepackage{amsmath,amssymb,amsthm,amsfonts}
\usepackage{geometry}
\geometry{margin=1in}
\usepackage{hyperref}

%-------------------------------------------------
%  独自マクロ定義（必要に応じて）
%-------------------------------------------------
\newcommand{\mweq}{\mathrel{\overset{\mathrm{mw}}{=}}}
\newcommand{\dfix}{D_{\mathrm{fix}0}}
\newcommand{\meta}{\Uparrow}

%-------------------------------------------------
%  定理環境の設定
%-------------------------------------------------
\theoremstyle{definition}
\newtheorem{principle}{原理}[section]

\begin{document}

\title{抽象原理：フラクタル性と閉じた抽象の限界}
\author{Masaomi Chiba}
\date{2026-04-22}
\maketitle

\begin{abstract}
本稿は、界論（Kairon）における「抽象化（Abstraction）」の数学的・トポロジカルな本質を定義する。あらゆる現象が無限のフラクタル構造を持つことを公理とし、理論的抽象とは、このフラクタルから根を切り離し、本来閉じない現実を強制的に閉包する操作であると位置づける。この「切り離し」こそが、局所的な排中律の成立と不完全性定理の根源であり、カオスは観測者の解像度限界によって生じるトポロジカルな散逸であることを示す。
\end{abstract}

\section{抽象原理}

\begin{principle}[現象のフラクタル性]
あらゆる現象はフラクタル性をもつ。
\end{principle}

\begin{principle}[抽象による根の切断]
あらゆる理論の抽象は、フラクタルから根を切り離すことである。
\end{principle}

\begin{principle}[強制的な閉包]
あらゆる理論の抽象は、本来は閉じない現実を閉じるものである。
\end{principle}

\begin{principle}[部分と全体の不可分性]
フラクタルは部分に切り出すことができない。部分が自己言及性を含むことは、部分が全体とは切り離せないことを意味する。
\end{principle}

\begin{principle}[閉じた抽象の整合性喪失]
切り出された抽象は、閉じた部分が原因となり全体との整合性を損なう。
\end{principle}

\begin{principle}[不完全性定理と排中律の発生]
切り出された部分には不完全性定理が生じ、排中律が成立する。
\end{principle}

\begin{principle}[カオスの観測論的成立]
カオスはフラクタル性を持つが、フラクタル性を全て観測できないことにより成立する。
\end{principle}

\end{document}
